% TOP_PROBLEM: {"id": 3700006, "title": "Completely invariant Fatou components", "category_id": 11, "category": {"id": 11, "name": "dynamical_systems", "display_name": "Dynamical Systems", "description": "Problems about long-term behavior of deterministic systems, Hamiltonian dynamics, and periodic orbits.", "slug": "dynamical-systems", "order_index": 11, "created_at": "2026-07-31T15:26:25.671Z"}, "status": "partially_solved", "problem_number": "AMR-036-0006", "set_id": 13, "statusQualification": "Uses the corrected problem following the documented gap in Baker's historical proof; no obsolete at-most-one theorem is asserted."}
% FIRST_POSTED: 2026-10-01
% ENTRY_KIND: statement_only
% UnsolvedMath ID 3700006; source catalogue code AMR-036-0006
% !TeX program = lualatex
% Definition and sources only; this file is not an LLM solution attempt.
\documentclass[11pt,a4paper]{article}
\usepackage[margin=25mm]{geometry}
\usepackage{fontspec}
\setmainfont{lmroman10-regular.otf}[BoldFont=lmroman10-bold.otf,
 ItalicFont=lmroman10-italic.otf,BoldItalicFont=lmroman10-bolditalic.otf]
\usepackage{amsmath,amssymb,mathtools}
\usepackage{unicode-math}
\setmathfont{latinmodern-math.otf}
\AtBeginDocument{\let\setminus\smallsetminus}
\usepackage{xurl,hyperref}
\hypersetup{colorlinks=true,urlcolor=blue}
\setcounter{secnumdepth}{0}
\setlength{\parindent}{0pt}
\setlength{\parskip}{5pt}
\setlength{\emergencystretch}{3em}
\begin{document}
\section*{Completely invariant Fatou components}\label{3700006}
{\small UnsolvedMath ID 3700006 · AMR-036-0006 · Dynamical Systems · AMR Open Problem Lists}
\subsection{Definitions and mathematical statement}
Let $f:\mathbb C\to\mathbb C$ be a transcendental entire
function, meaning an entire function that is not a polynomial.
Its Fatou set $F(f)$ consists of those points with a neighborhood
on which the iterates $\{f^m:m\geq0\}$ form a normal family
in the spherical metric. A Fatou component is a maximal
connected open subset of $F(f)$.

A Fatou component $U$ is completely invariant if
\[
                      f^{-1}(U)=U.
\]
Thus $f(U)\subseteq U$, and every point whose image lies in
$U$ already belongs to $U$. This full-preimage condition is
stronger than mere forward invariance of the component.
Surjectivity onto values omitted by $f$ is not required.

What is the largest possible number of completely invariant
Fatou components of a transcendental entire function? In
particular, is that number always at most one, or can there
exist two disjoint Fatou components $U,V$ with
$f^{-1}(U)=U$ and $f^{-1}(V)=V$?
The target covers all transcendental entire functions,
without assuming finite order, finitely many singular values,
or any special topology of their Fatou components.

Any negative example to the at-most-one assertion must use
actual components of the dynamical Fatou set. Merely finding
disjoint plane domains whose preimages are connected does
not give such an example. Conversely, an upper bound must
apply uniformly to the full class of entire functions.
\subsection{Short English statement}
Can a transcendental entire function have two completely invariant Fatou components?
\subsection{Sources}
\begin{itemize}
\item[S1] ULAM.AI and the UnsolvedMath Contributors, supplied problems.json, record 3700006 (AMR-036-0006), AMR Open Problem Lists. Catalogue identity and source transcription; CC BY 4.0.
\url{https://huggingface.co/datasets/ulamai/UnsolvedMath}.
\item[S2] Eremenko - My favorite unsolved problems, Completely invariant domains, Problem 2. Original source indexed by AMR-036-0006.
\url{https://www.math.purdue.edu/~eremenko/uns1.html}.
\item[S3] A. Eremenko, Two questions on completely invariant domains, Question 2. Author-maintained primary problem note.
\url{https://www.math.purdue.edu/~eremenko/dvi/conjectures.pdf}.
\end{itemize}
\end{document}
